%======================================================================
\section{Analytic details}\label{sec:analysis}
%======================================================================
The proposal reduces to analytic statements of four kinds:
\begin{enumerate}[label=(\alph*)]
\item an adiabatic limit with a knot transverse to the neck;
\item positivity and admissibility of the charges in that limit;
\item an intrinsic, neck-free description of the charges near $K$;
\item reducibles, compactness and invariance with the new weights.
\end{enumerate}
We take them in turn and say what is standard, what is new, and where the difficulty lies.

\subsection{The adiabatic limit with a transverse knot}
\emph{Monopoles.}  Use the metric $dt^2+ds^2+\varepsilon^2g_\Sigma$ on $\R\times N_T$, with an area-form perturbation of strength
$\tau$ satisfying the Bradlow condition for degree-$g$ vortices.  As $\varepsilon\to0$, solutions on the neck converge, away from
a finite bubbling set, to holomorphic maps $\R\times[-T,T]\to\mathrm{Sym}^g\Sigma$, the space of degree-$g$ $\tau$-vortices.  This is
Salamon's analysis for mapping tori \cite{Salamon-SW}, transplanted to a neck with Lagrangian ends.  What is new is
\emph{matching at the ends of the neck}.  One must show that the boundary values at $s=\pm T$ of finite-energy solutions
on the handlebody pieces $U_\alpha^\circ,U_\beta^\circ$ (cylindrical ends $[0,\infty)\times\Sigma$) converge to $\mathbb T_\alpha,\mathbb T_\beta$.  The model
statement is: \emph{the vortex boundary values of $\tau$-perturbed SW solutions on a handlebody with product end are the
divisors supported on the $\alpha$-curves, for a suitable choice of $\alpha$-curves determined by the metric}.  This is
the SW counterpart of ``flat connections on a handlebody form $L_\alpha$''.  It is the hardest single input and amounts
to the Atiyah--Floer conjecture for Seiberg--Witten.  The knot adds nothing to it: away from $s=\pm T$ the knot is the
pair of lines $\{w,z\}\times[-T,T]$, which are complex curves for the product structure.

\emph{Instantons.}  The analysis here is more advanced: Dostoglou--Salamon for mapping tori \cite{DoSa}, Salamon--Wehrheim for
Lagrangian boundary conditions \cite{SW-lag}, Daemi--Fukaya--Lipyanskiy for the $SO(3)$ Atiyah--Floer conjecture
\cite{DFL}.  The new ingredient is the determinant-line section along the strips, which is a smooth function of the
connection on the slices and so passes to the limit with the connection.  In its four-dimensional form, $c_z$ is the relative index of the Dirac
operator on $\R\times Y$ coupled to $\mathbb A\otimes\lambda_z$, with APS conditions at the ends.  It is therefore defined, additive and
locally constant for \emph{every} configuration.  This is cleaner than the monopole case: no admissibility is needed, and
only positivity remains, via Problem~\ref{prob:avoid}.  The analytic input is a gluing/adiabatic formula for this relative
index under neck-stretching, identifying it with the relative degree of $\mathcal D_z$ on the limiting strip plus bubble
contributions.  This is a families-index computation of Bismut--Freed/Dai type.

\subsection{Positivity in the limit}
\begin{lemma}[stability of non-negative zero counts]\label{lem:posvortex}
Let $\Omega=\R\times[-T,T]$, and let $s_\varepsilon$ be sections of a trivialised line bundle over $\Omega$ that converge in $C^0_{\rm loc}$ to a
holomorphic section $s_0$, uniformly on $\R\times\{\pm T\}$ and on $\{|t|\ge R_0\}\times[-T,T]$, where $|s_0|\ge c>0$.  Then the
signed zero count of $s_\varepsilon$ equals that of $s_0$ for $\varepsilon$ small, and it is $\ge0$.
\end{lemma}
\begin{proof}
On $\partial([-R_0,R_0]\times[-T,T])$ we have $|s_\varepsilon-s_0|<c\le|s_0|$ for small $\varepsilon$, so the winding numbers agree (Rouché).  The
zeros of $s_0$ are isolated and positive.
\end{proof}
The analytic work is in the hypotheses.  The first, $C^0$ convergence of the vortex component $\phi^-_z$ \emph{at the
specific point} $z$, needs elliptic estimates along $\Sigma$ that are uniform in $\varepsilon$.  These are available away from
bubbling points.  At a bubbling point lying on $\R\times\gamma^T_z$ one needs the bubble's own zero count.  In
$\mathrm{Sym}^g\Sigma$, sphere bubbles meet $V_w$ and $V_z$ equally and positively, and boundary degenerations carry $UV$.  The
gauge-theoretic statement to prove is that the energy-concentration profiles in the $\varepsilon\to0$ limit are vortex bubbles
of this type, and never configurations in which $\phi^-_z$ winds negatively.  The second hypothesis, uniformity in
$t$, is exponential convergence to critical points with constants independent of $\varepsilon$.  Since the
neck-stretched critical points correspond to transverse intersections of $\mathbb T_\alpha,\mathbb T_\beta$, this is a spectral-gap statement
for the Hessian, uniform in $\varepsilon$, of Dostoglou--Salamon type.

With Lemma~\ref{lem:bootstrap} these give Conjecture~\ref{conj:Mpos} for $T\ge T(E)$.  The energy window is essential:
there is no reason for positivity to hold uniformly over all energies at a fixed neck length.  The invariant is defined
as a direct limit (\S\ref{subsec:inv}).

\subsection{The intrinsic picture near $K$: critical holonomy and creases}\label{subsec:edge}
The neck is a scaffold.  For an invariant of $(Y,K)$ one wants to know what the charges are near $K$ when no splitting
is present.  The neck description says that along $\gamma_w$ the relevant component is $E_+(\dot K)$, and along $\gamma_z$ it is
$E_-(\dot K)$, switching at $p_0,p_3$.  Here is a local model in which this switching is built into the equations.

\paragraph{Aharonov--Bohm holonomy.}  Couple the spin$^c$ connection to a flat $U(1)$ connection on $Y\setminus K$ with
meridional holonomy $e^{2\pi i\beta}$.  For $U(1)$ this is the limit of an \emph{exact} perturbation, since the Thom form of a
null-homologous $K$ is exact.  So $\HM$ is unchanged for smooth smearings, which is \ref{Q:base}.  In the singular limit,
transversally to $K$ the Dirac operator is the Aharonov--Bohm operator on $\C$ with flux $\beta\in(0,1)$.  Its local
solutions near $r=0$ have leading terms
\[
\Phi\sim c_+\,r^{-\beta}\,\mathbf e_+ +c_-\,r^{\beta-1}\,\mathbf e_-+O(r^{\min(1-\beta,\beta)}),\qquad \mathbf e_\pm\in E_\pm(\dot K),
\]
and both singular terms are in $L^2$.  The operator therefore has a family of closed extensions, parametrised by
Lagrangians in the space of edge data $(c_+,c_-)\in\Gamma(K;E_+\oplus E_-)$, with the symplectic form coming from Green's formula,
which pairs $c_+$ with $c_-$.  Two extensions are distinguished: $\mathcal E_w=\{c_-=0\}$, which tends to the smooth theory as
$\beta\to0$, and $\mathcal E_z=\{c_+=0\}$, which tends to the smooth theory after the singular gauge transformation $e^{i\theta}$ as
$\beta\to1$.  At the \emph{critical holonomy} $\beta=\tfrac12$ both singular terms are $r^{-1/2}$, and the two extensions are
exchanged by the reflection $\theta\mapsto-\theta$.

\paragraph{Creases.}  The proposal is to use $\beta=\tfrac12$ with the boundary condition $\mathcal E_w$ along $\gamma_w$ and $\mathcal E_z$ along
$\gamma_z$.  The condition switches at the two \emph{creases} $p_0,p_3$, and the surviving edge coefficient is the section
$c=(c_+|_{\gamma_w},c_-|_{\gamma_z})$.  This is precisely the neck description, now intrinsic, and the charges become
\[
n_w=\#\,c_+^{-1}(0)\cap(\R\times\gamma_w),\qquad n_z=\#\,c_-^{-1}(0)\cap(\R\times\gamma_z).
\]
Their difference is a topological count for the glued line bundle over $\R\times K$, which gives \ref{M:alex}.  Positivity on
both arcs is what prevents zeros from crossing a crease: a positive zero of $c_+$ on the $\gamma_w$ side would arrive as a
negative zero of $c_-$ on the $\gamma_z$ side.  So the creases behave like the Lagrangian walls of the cylindrical
reformulation of Heegaard Floer.  On a cobordism with surface $F$, the boundary condition can switch along curves on
$F$.  These \emph{crease curves are Zemke's dividing arcs} \ref{M:cob}, and the index contributions of the
$\mathcal E_w$- and $\mathcal E_z$-regions should reproduce the $\tilde\chi(F_w)$ and $\tilde\chi(F_z)$ terms of his grading formula.

\paragraph{What fails, and what must be proved.}
\begin{enumerate}[label=(\roman*),leftmargin=2em]
\item \emph{Edge calculus.}  The linearised operators are edge operators in Mazzeo's sense \cite{Mazzeo} along $K$ (in $3$d) and
$\R\times K$ (in $4$d), with indicial roots $\pm\frac12+\Z$ for the Dirac part at $\beta=\frac12$.  Fredholmness on weighted spaces with
the mixed boundary condition, and an index formula, are standard in principle.  The crease is a codimension-$3$
stratum in $4$d where the boundary condition jumps.  Its local model is a two-dimensional problem with corner-type
indicial roots.  It must be shown that the index jumps by an amount independent of the configuration, which is the
analogue of the fixed Maslov contribution of a basepoint region.
\item \emph{$C^0$ bound.}  The standard bound $|\Phi|^2\le\max(-s_g,0)$ fails, since $|\Phi|\sim r^{-1/2}$ near $K$.  Replace it by a
weighted maximum principle for $r|\Phi|^2$, using the Weitzenböck formula on $Y\setminus K$ together with the boundary term from
Green's formula.  That boundary term is $\mathrm{Re}\int_K\langle c_+,c_-\rangle$ up to constants, and it \emph{vanishes} for $\mathcal E_w$ and for
$\mathcal E_z$.
\item \emph{Edge bubbling.}  The SW equations are scale-invariant on $\R^4$, and only the $C^0$ bound prevents bubbling.
Near $K$, rescaled limits are finite-energy solutions on $\R^4$ with Aharonov--Bohm flux $\frac12$ along $\R^2$.  For a
one-sided condition the Weitzenböck identity $\int|\nabla\Phi|^2+\frac14|\Phi|^4=(\text{edge term})=0$ makes them trivial.  At a crease
the edge term need not vanish, so nontrivial \emph{crease bubbles} may exist.  The prediction is that they exist and
carry weight $UV$, cancel in pairs, and are the gauge-theoretic boundary degenerations of
Proposition~\ref{prop:d2}.  Classifying finite-energy solutions of this local model is a concrete,
finite-dimensional-symmetry problem.  By $\mathrm{SO}(2)\times\mathrm{SO}(2)$-invariance it reduces to ODE.
\item \emph{Positivity without a neck.}  Unlike in the neck limit, the edge coefficients satisfy no local
$\bar\partial$-equation: leading edge data are free boundary data.  Intrinsic positivity of $n_w,n_z$ therefore requires a
global argument, for instance a Kähler-type model metric on a tubular neighbourhood of $\R\times K$ combined with a
large perturbation localised there, in Taubes' style.  We regard the neck construction as primary and this as its
eventual intrinsic replacement.
\end{enumerate}

\paragraph{Bordered consistency check.}  If one cuts along $T^2=\partial\nu K$ instead of $\Sigma$, the boundary phase space is the
dual torus $\Tdual=H^1(T^2;\R/\Z)$ of flat $U(1)$ connections.  The flat-torus Dirac operator has kernel at exactly one point
$\theta_0\in\Tdual$ for a fixed spin structure, and there the kernel is $\C\oplus\C$, one constant spinor of each chirality.  Three
facts can be checked directly:
\begin{enumerate}[label=(\arabic*)]
\item the solid-torus Lagrangian avoids $\theta_0$;
\item the reducible locus of the knot complement, $\{\mathrm{hol}_\lambda=1\}$, avoids $\theta_0$;
\item the Aharonov--Bohm filling at $\beta=\frac12$ passes through $\theta_0$.
\end{enumerate}
This matches the Hanselman--Rasmussen--Watson picture.  Fillings miss the marked point, $\widehat{\HFK}$ is computed by a meridian
through it, the essential component of the curve is the reducible line, and the two chiral kernels at $\theta_0$ are the two
sides of the marking that carry $w$ and $z$.  It is an independent check that critical holonomy is the right local
model.

\subsection{Reducibles, compactness and invariance}\label{subsec:inv}
\begin{itemize}[leftmargin=2em]
\item \emph{Blow-up.}  In Kronheimer--Mrowka's $\sigma$-model the charges are defined with the normalised spinor $\phi$, so they
make sense on reducible trajectories.  For a flow from eigenvector $\psi_\lambda$ to eigenvector $\psi_\mu$, $n_z$ is the difference of the
windings of $\pi_{E_-}\psi_\lambda$ and $\pi_{E_-}\psi_\mu$ along $\gamma_z^T$, or of the edge coefficients in the intrinsic model.  Genericity of
eigenvectors, meaning non-vanishing of a section of a real rank-$2$ bundle along an arc, is a codimension-one condition
in $1$-parameter families.  It must be arranged by perturbations that preserve the structure near $K$.
\item \emph{Energy windows and direct limits.}  Take $C_{T,E}$ to be the subquotient generated in energy $\le E$.  For
$T'\ge T\ge T(E')$ the continuation maps $C_{T,E}\to C_{T',E'}$ are $V$-weighted and positive, and the invariant is
$\varinjlim_E H_*(C_{T(E),E})$.  This copies Taubes' passage from $\HM$ to ECH via action-filtered complexes.  The point to
check is that continuation trajectories also satisfy positivity once $T$ is large relative to the energy.  That is the
1-parameter version of Lemma~\ref{lem:posvortex}.
\item \emph{Splitting independence.}  Stabilisations and handleslides away from $w,z$ are realised by stretching along
the extra surfaces of a Heegaard triple.  The corresponding triangle counts are $V$-weighted by the same charges on the
triangle's neck strips.  Moves across $z$ are exactly those along which admissibility fails, as in Heegaard Floer.
\item \emph{Orientations.}  Work over $\F$ first.  Over $\Z$, Kronheimer--Mrowka's orientation of the determinant line is
unaffected by the weights, but crease bubbles (if they exist) must be shown to cancel with signs.  In Heegaard Floer
the two boundary degenerations do cancel with signs.
\end{itemize}
